\documentclass[a4paper,11pt]{jsarticle}
\usepackage[margin=1cm]{geometry}
\usepackage{amsmath, amssymb}
\usepackage{multicol}

% ページ番号を消す設定
\pagestyle{empty}

% ページを真ん中で分ける線の設定
\setlength{\columnsep}{1cm}
\setlength{\columnseprule}{0.4pt}

\newcommand{\testheader}{
\begin{center}
{\large 代数概論・同演習 期末試験(2024/1/26)}\\[1em]
\underline{\hspace{1cm}}年\underline{\hspace{1cm}}組\underline{\hspace{1cm}}番\hspace{0.5cm}名前\underline{\hspace{6cm}}
\end{center}
}

\begin{document}
% ==========================================
% ヘッダー (1ページ目の全体の上)
% ==========================================
\noindent
%\newcommand{\testheader}{
%\begin{center}
%{\large 代数概論 期末試験(2024/1/26)}\\[1em]
%\underline{\hspace{1cm}}年\underline{\hspace{1cm}}組\underline{\hspace{1cm}}番\hspace{0.5cm}名前\underline{\hspace{6cm}}
%\end{center}
%\begin{multicols*}{2}
%}

% ==========================================
% 1ページ目 左側 (領域1)
% ==========================================
\testheader
\begin{multicols*}{2}

\noindent
\textbf{問1.}\quad 次の用語や定理について説明せよ．\\[1em]
(1)[3点] 最小多項式
\vspace{4cm}

\noindent
(2)[3点] 長さ $l$ の巡回置換
\vspace{4cm}

\noindent
(3)[3点] 代数的数
\vspace{4cm}

\noindent
(4)[3点] アイゼンシュタインの既約判定法
\vspace{4cm}

\noindent
(5)[3点] 対称群の元 $\sigma$ の位数
\vspace{4cm}

\vfill\null
\columnbreak
% ==========================================
% 1ページ目 右側 (領域2)
% ==========================================
\noindent
\textbf{問2.}[7点] 整数係数モニック多項式 $f(x)$ について，有理数 $r$ が $f(r)=0$ を満たすならば，$r$ は整数であることを示せ．
\vspace{15cm}

\noindent
\textbf{問3.}[4点] 変数 $x_1, x_2, x_3$ からなる次の7つの単項式を，辞書式順序で並べ替えよ．\\[1em]
$x_1^2 x_2, \quad x_1x_2^3, \quad x_1 x_2 x_3, \quad x_2x_3, \quad x_1 x_3^2, \quad x_1^3, \quad x_2^2 x_3$
\vspace{7cm}

\vfill\null

\newpage
% ==========================================
% 2ページ目 左側 (領域3)
% ==========================================
\noindent
\textbf{問4.}\quad 対称群 $\mathfrak{S}_9$ の元 $\sigma$ を次のように定める．
$$\sigma = \begin{pmatrix} 1 & 2 & 3 & 4 & 5 & 6 & 7 & 8 & 9 \\ 3 & 5 & 1 & 6 & 2 & 4 & 9 & 7 & 8 \end{pmatrix}$$
次の問いに答えよ．\\
(1)[4点] $\sigma$ を巡回置換の積として表せ．
\vspace{4cm}

\noindent
(2)[4点] $\sigma$ の位数を求めよ．
\vspace{4cm}

\noindent
(3)[4点] $\sigma$ を互換の積として表せ．
\vspace{4cm}

\noindent
(4)[4点] $\sigma^2$を求めよ．
\vspace{4cm}

\noindent
(5)[4点] 逆元 $\sigma^{-1}$ を求めよ．
\vspace{4cm}

\vfill\null
\columnbreak
% ==========================================
% 2ページ目 右側 (領域4)
% ==========================================
\noindent
\textbf{問5.}[8点] 多項式 $f(x) = 1$ と $g(x) = b_m x^m + b_{m-1}x^{m-1} + \dots + b_0$ について，終結式$R_{1,m}(f,g)$を計算せよ．
\vspace{22cm}

\vfill\null
\end{multicols*}
\newpage
% ==========================================
% 3ページ目 左側 (領域5)
% ==========================================
\testheader
\begin{multicols*}{2}
\noindent
\textbf{問6.}\quad 多項式 $f(x) = a_n x^n + \dots + a_1 x + a_0$ と $g(x) = b_m x^m + \dots + b_1 x + b_0$ について，次の問いに答えよ．\\[0.5em]
(1)[5点] 積 $f(x)g(x)$ を$a_n$と$b_n$を用いて表せ．
\vspace{7cm}

\noindent
(2)[7点] $f(x)$と$g(x)$が原始多項式のとき，積$f(x)g(x)$も原始多項式となることを示せ．
\vspace{14.5cm}

\vfill\null
\columnbreak
% ==========================================
% 3ページ目 右側 (領域6) 
% ==========================================
\noindent
\textbf{問7.}\quad 次の多項式が有理数体 $\mathbb{Q}$ 上既約であることを示せ．\\
(1)[5点] $x^4 - 14x^3 + 21x^2 - 35x + 14$
\vspace{7cm}

\noindent
(2)[5点] $x^4 + x^3 + x^2 + x + 1$
\vspace{7cm}

\noindent
(3)[5点] $x^3 - x - 1$
\vspace{7cm}

\vfill\null
\newpage
% ==========================================
% 4ページ目 左側 (領域7)
% ==========================================
\noindent
\textbf{問8.}\quad 変数$x,y,z,w$について，次の問いに答えよ．\\
(1)[4点] 4次基本対称式 $e_1, e_2, e_3, e_4$ をそれぞれ求めよ．
\vspace{5cm}

\noindent
(2)[7点] 次の対称式を，基本対称式 $e_1, e_2, e_3, e_4$ を用いて表せ．
$$ f(x,y,z,w)=x^3 + y^3 + z^3 + w^3 - 4xyzw$$ 
\vspace{10cm}

\vfill\null
\columnbreak
% ==========================================
% 4ページ目 右側 (領域8)
% ==========================================
\noindent
\textbf{問9.}[8点] 次の式は交代式である．因数分解せよ．
$$ f(x,y,z)=x^4(y-z) + y^4(z-x) + z^4(x-y) $$
\vspace{22cm}

\vfill\null
\end{multicols*}
\end{document}